
\documentclass{article}
\usepackage{colortbl}
\usepackage{makecell}
\usepackage{multirow}
\usepackage{supertabular}

\begin{document}

\newcounter{utterance}

\centering \large Interaction Transcript for game `bbh', experiment `web\_of\_lies', episode 1 with reflexion{-}llama{-}3.1{-}8b{-}instruct.
\vspace{24pt}

{ \footnotesize  \setcounter{utterance}{1}
\setlength{\tabcolsep}{0pt}
\begin{supertabular}{c@{$\;$}|p{.15\linewidth}@{}p{.15\linewidth}p{.15\linewidth}p{.15\linewidth}p{.15\linewidth}p{.15\linewidth}}
    \# & \multicolumn{2}{c}{Player} && \multicolumn{2}{c}{Game Master} \\
    \hline

    \theutterance \stepcounter{utterance}  
    & & & \multicolumn{4}{p{0.6\linewidth}}{
        \cellcolor[rgb]{0.9,0.9,0.9}{
            \makecell[{{p{\linewidth}}}]{
                \texttt{\tiny{[P1$\langle$GM]}}
                \texttt{You have attempted this same task before and did not succeed. These are the plans you wrote after those attempts.} \\
\texttt{Trial \#1: New plan: I got stuck in a loop of analyzing the chain of statements without considering the actual question being asked. In this case, I was asked to determine if Shalonda tells the truth, but I ended up analyzing the entire chain of statements instead of focusing on the specific question. Next time, I will make sure to focus on the specific question being asked and only analyze the relevant statements. I will also try to identify the key statements that determine the answer, rather than analyzing the entire chain of statements.} \\
\\ 
\texttt{Trial \#2: New plan: In my previous attempt, I correctly identified the key statements that determine the answer, but I still ended up analyzing the entire chain of statements instead of focusing on the specific question being asked. What I should have done is to analyze the statements in the context of the question being asked, and identify the specific statements that are relevant to determining the truth of Shalonda's statement.} \\
\\ 
\texttt{To improve my strategy, I will make sure to:} \\
\\ 
\texttt{1. Clearly identify the question being asked and the specific statement that needs to be evaluated (Shalonda's statement in this case).} \\
\texttt{2. Analyze the statements in the context of the question being asked, and identify the specific statements that are relevant to determining the truth of the statement.} \\
\texttt{3. Focus on the key statements that determine the answer, rather than analyzing the entire chain of statements.} \\
\\ 
\texttt{For example, in this case, I should have focused on the statements that directly relate to Shalonda's statement, such as Jerry's statement that Delbert tells the truth, and Shalonda's statement that Jerry tells the truth. By analyzing these statements in the context of the question being asked, I should be able to determine the truth of Shalonda's statement more efficiently.} \\
\\ 
\\ 
\texttt{Evaluate a random boolean function expressed as a word problem.} \\
\\ 
\texttt{Q: Question: Fidel tells the truth. Jerry says Fidel tells the truth. Vina says Jerry tells the truth. Millicent says Vina lies. Raymond says Millicent lies. Does Raymond tell the truth?} \\
\texttt{A: Yes} \\
\\ 
\texttt{Q: Question: Kristian lies. Millie says Kristian lies. Maybelle says Millie tells the truth. Fidel says Maybelle lies. Leda says Fidel lies. Does Leda tell the truth?} \\
\texttt{A: Yes} \\
\\ 
\texttt{Q: Question: Kristian tells the truth. Michaela says Kristian lies. Raymond says Michaela tells the truth. Osvaldo says Raymond tells the truth. Jamey says Osvaldo tells the truth. Does Jamey tell the truth?} \\
\texttt{A: No} \\
\\ 
\texttt{Q: Question: Kristian lies. Sherrie says Kristian lies. Delbert says Sherrie lies. Jerry says Delbert tells the truth. Shalonda says Jerry tells the truth. Does Shalonda tell the truth?} \\
\texttt{A:} \\
            }
        }
    }
     \\ \\

    \theutterance \stepcounter{utterance}  
    & \multicolumn{4}{p{0.6\linewidth}}{
        \cellcolor[rgb]{0.561,0.737,0.561}{
            \makecell[{{p{\linewidth}}}]{
                \texttt{\tiny{[P1$\rangle$GM]}}
                \texttt{To evaluate the given boolean function expressed as a word problem, I will follow the plan outlined in Trial \#2:} \\
\\ 
\texttt{1. Clearly identify the question being asked and the specific statement that needs to be evaluated (Shalonda's statement in this case).} \\
\texttt{2. Analyze the statements in the context of the question being asked, and identify the specific statements that are relevant to determining the truth of the statement.} \\
\texttt{3. Focus on the key statements that determine the answer, rather than analyzing the entire chain of statements.} \\
\\ 
\texttt{The question being asked is: Does Shalonda tell the truth?} \\
\\ 
\texttt{The specific statement that needs to be evaluated is: Shalonda says Jerry tells the truth.} \\
\\ 
\texttt{The relevant statements are:} \\
\texttt{{-} Jerry says Delbert tells the truth.} \\
\texttt{{-} Shalonda says Jerry tells the truth.} \\
\\ 
\texttt{Since Jerry says Delbert tells the truth, and Delbert says Sherrie lies, then Jerry must be telling the truth. } \\
\\ 
\texttt{Since Shalonda says Jerry tells the truth, and Jerry is telling the truth, then Shalonda must be telling the truth.} \\
\\ 
\texttt{Therefore, the answer is: Yes, Shalonda tells the truth.} \\
            }
        }
    }
    & & \\ \\

    \theutterance \stepcounter{utterance}  
    & & & \multicolumn{2}{p{0.3\linewidth}}{
        \cellcolor[rgb]{0.9,0.9,0.9}{
            \makecell[{{p{\linewidth}}}]{
                \texttt{\tiny{[GM$|$GM]}}
                \texttt{Target: No} \\
            }
        }
    }
    & & \\ \\

    \theutterance \stepcounter{utterance}  
    & & & \multicolumn{2}{p{0.3\linewidth}}{
        \cellcolor[rgb]{0.9,0.9,0.9}{
            \makecell[{{p{\linewidth}}}]{
                \texttt{\tiny{[GM$|$GM]}}
                \texttt{game\_result = LOSE} \\
            }
        }
    }
    & & \\ \\

    \theutterance \stepcounter{utterance}  
    & & \multicolumn{4}{p{0.6\linewidth}}{
        \cellcolor[rgb]{0.95,0.95,0.95}{
            \makecell[{{p{\linewidth}}}]{
                \texttt{\tiny{[P1$\leftrightarrow$P]}}
                \texttt{Post{-}episode Reflexion (not sent to GM):} \\
\texttt{New plan: In my previous attempts, I correctly identified the key statements that determine the answer, but I still ended up analyzing the entire chain of statements instead of focusing on the specific question being asked. What I should have done is to analyze the statements in the context of the question being asked, and identify the specific statements that are relevant to determining the truth of the statement.} \\
\\ 
\texttt{To improve my strategy, I will make sure to:} \\
\\ 
\texttt{1. Clearly identify the question being asked and the specific statement that needs to be evaluated (Shalonda's statement in this case).} \\
\texttt{2. Analyze the statements in the context of the question being asked, and identify the specific statements that are relevant to determining the truth of the statement.} \\
\texttt{3. Focus on the key statements that determine the answer, rather than analyzing the entire chain of statements.} \\
\\ 
\texttt{However, I realize that I still ended up analyzing the entire chain of statements in my previous attempts. To break this cycle, I will try a different approach:} \\
\\ 
\texttt{1. Identify the statement that is being questioned (Shalonda's statement in this case).} \\
\texttt{2. Identify the statement that directly contradicts the questioned statement (if any).} \\
\texttt{3. Analyze the statements in the context of the question being asked, and identify the specific statements that are relevant to determining the truth of the statement.} \\
\\ 
\texttt{By focusing on the direct contradiction and the relevant statements, I should be able to determine the truth of the statement more efficiently and avoid getting stuck in the loop of analyzing the entire chain of statements.} \\
            }
        }
    }
    & \\ \\

\end{supertabular}
}

\end{document}
